\section{Introduction}
An automatic post-editor must decide both what to correct and what to leave alone. A fluent rewrite can repair an omission, but it can also alter a name, number, or source relation that the machine translation (MT) draft already expressed correctly. When drafts are strong, mean quality changes therefore give an incomplete account of editing behaviour. LLMs can improve neural MT output, yet they also produce hallucinated edits~\cite{raunak2023}.

We study this tension for English--Tagalog and English--Spanish in both directions. A fixed NLLB-200 3.3B system~\cite{nllb2022} produces sentence-split drafts, and two instruction-following editors, open-weight Llama 3 8B Instruct and commercial GPT-4o Mini, receive the source and the draft. The two language pairs provide a limited resource contrast; language, direction, domain, and baseline quality remain intertwined, so we report pair and direction differences descriptively rather than as causal resource effects. We ask:
\textbf{RQ1:} How do editing gains and harms differ between the Tagalog and Spanish pairs, conditional on baseline quality?
\textbf{RQ2:} How do the three prompting strategies compare within each model and direction?
\textbf{RQ3:} When do editors change already strong drafts and degrade their measured quality or source fidelity?
\textbf{RQ4:} How do effects differ between the two directions of a language pair?

This is an empirical evaluation, not a new post-editing architecture. Anchoring an LLM editor on a frozen NMT draft and selecting in-context exemplars by semantic similarity follow established prior work~\cite{ahmadnia2026icai,vilar2023}, and we do not claim either as a contribution. Our contributions are:
\begin{itemize}
\item a held-out evaluation in a new empirical setting (four Tagalog/Spanish directions) against a stronger, sentence-split NLLB-200 baseline, with the protocol frozen on development data before the test partition was opened;
\item a two-editor comparison (open-weight and commercial) across three prompting strategies, with paired confirmatory tests and Holm multiplicity control;
\item explicit failure accounting: every planned output is retained, and output-contract failures are reported as outcomes rather than discarded;
\item an over-editing and preservation analysis of strong drafts that separates metric-defined harm from heuristic preservation screens.
\end{itemize}
